\chapter{Conclusiones y trabajo futuro}

El sistema aprende una ordenación transversal fuera de muestra: su Rank-IC medio es 0,1004, supera
los baselines deterministas y se mantiene positivo tras placebos, permutación, bootstrap, exclusión
de eras, semillas y neutralización por estilo. El meta-agente aporta valor frente al promedio
equiponderado porque aprende a reponderar especialistas a partir de cohortes cerradas.

La traducción económica es positiva en la ventana de selección y en los dos años reservados. Sin
embargo, la conclusión no es una promesa de inversión: el Deflated Sharpe no supera el umbral, la
muestra efectiva es pequeña y la confirmación de Rank-IC en la era reservada aún no tiene potencia.
La contribución defendible del TFM es metodológica: demuestra cómo separar señal, selección,
robustez y cartera en un sistema de aprendizaje financiero reproducible.

La prioridad futura no es añadir complejidad al modelo, sino elevar el coeficiente de transferencia:
ampliar el número de posiciones, reducir rotación, estudiar sizing por convicción y probar costes
dependientes de liquidez. También será necesario esperar el cierre de nuevas cohortes reservadas y
replicar el protocolo en otros universos y regímenes.

\section{Respuesta a la pregunta de investigación}

La respuesta es afirmativa con condiciones. Dentro del universo, periodo y protocolo definidos, los
cinco agentes y su meta producen una ordenación transversal con Rank-IC positivo, superior a los
baselines incluidos y estable frente a contrastes adversariales. El meta parte de pesos iguales y
aprende de cohortes cerradas, por lo que su comportamiento no se reduce a fijar de antemano el
especialista ganador.

La respuesta económica es positiva pero más acotada. La cartera supera a SPY en selección y en los
dos años reservados neta de los costes modelados. Sin embargo, el Deflated Sharpe y la muestra
efectiva impiden transformar esa observación en una promesa de superioridad ajustada por riesgo. La
formulación defendible es que el ranking tiene evidencia predictiva robusta y que su implementación
es económicamente compatible con esa señal bajo restricciones declaradas.

\section{Contribución y agenda}

La contribución principal es un protocolo auditable: separa fecha de publicación y cierre fiscal,
etiqueta futura y entrada de entrenamiento, selección predictiva y cartera, y evidencia de selección
y confirmación reservada. Los artefactos permiten reconstruir configuraciones, pesos, decisiones,
figuras y tablas. Esta disciplina es transferible a otros problemas de ML financiero incluso si el
ranking concreto no se replica.

El trabajo futuro debe priorizar la transferencia antes que añadir complejidad al modelo: estudiar
tamaño de cartera, sizing, costes y rebalanceo sin reabrir la selección predictiva. Después debe
reducirse la multiplicidad con un catálogo pre-registrado, esperar etiquetas reservadas y replicar
el protocolo en otros universos. Cada extensión debe preservar la observabilidad temporal que hace
interpretables los resultados actuales.

\section{Agenda ordenada por prioridad}

El barrido diagnóstico del Capítulo~\ref{chap:resultados} no sólo señala dónde está el cuello de
botella: indica en qué orden conviene atacarlo.

\begin{enumerate}[leftmargin=*,itemsep=0.4em]
  \item \textbf{Medir el mínimo de tenencia como variable de pleno derecho.} Es la intervención con
  mejor relación entre evidencia disponible y esfuerzo: el barrido ya sugiere que
  \texttt{half\_horizon} domina simultáneamente en exceso, IR y rotación. Convertirlo en una
  decisión medida --- y no heredada del valor recomendado --- exigiría declararlo antes de observar
  el resultado, para no incurrir en la selección por rentabilidad que el protocolo prohíbe.
  \item \textbf{Resolver el conflicto entre cadencia y rotación.} La cadencia mensual ganó la primera
  decisión con una ventaja clara, pero es también la causa estructural del 359\,\% de rotación:
  re-decidir cada mes sobre una señal a doce meses. Separar la frecuencia de \emph{evaluación} de la
  de \emph{ejecución} permitiría conservar las 117 cohortes sin pagar la rotación asociada.
  \item \textbf{Reducir el catálogo antes de volver a ejecutar.} El Deflated Sharpe penaliza por las
  66 configuraciones ensayadas. Un catálogo pre-registrado más estrecho, que fije de antemano las
  variables cuyo valor no se espera que cambie, elevaría la probabilidad deflactada sin necesidad de
  mejorar el modelo.
  \item \textbf{Esperar el cierre de cohortes reservadas.} La confirmación predictiva de 2025--2026
  es hoy indeterminada por falta de etiquetas cerradas, no por resultado adverso. Es la única
  limitación de esta lista que se resuelve sola con el paso del tiempo.
  \item \textbf{Investigar por qué domina \textit{risk}.} Que un solo agente supere al conjunto en
  las cuatro eras es el resultado menos explicado del trabajo. Puede ser una prima de baja
  volatilidad conocida, un artefacto del universo de gran capitalización o una interacción con el
  horizonte de doce meses; distinguirlo exige un diseño específico.
\end{enumerate}

Obsérvese que ninguna de las cinco prioridades consiste en usar un modelo más potente. El sistema
captura una cuarta parte de la señal que ya tiene; mejorar la señal antes de resolver la
transferencia sería optimizar el término equivocado.

\section{Qué aprendería quien repitiera este trabajo}

Algunas lecciones del desarrollo, documentadas en el Capítulo~\ref{chap:desarrollo}, son
independientes del dominio financiero y probablemente más reutilizables que el resultado empírico.

La primera es que declarar no es computar. Un catálogo puede describir variables que el código nunca
calcula, y el sistema seguirá produciendo resultados plausibles: la comprobación de que lo declarado
coincide con lo computado debe ser automática, porque no es observable a simple vista.

La segunda es que una puerta estadística puede fallar en la dirección contraria a la esperada. La
regla de no inferioridad de este proyecto penalizaba a los candidatos superiores porque su intervalo
era más ancho, y el síntoma no fue un error sino una regularidad sospechosa: nadie resultaba
elegible nunca.

La tercera es que la ausencia de excepción no es evidencia de corrección. Un modelo que devuelve una
constante, una calibración que colapsa, una política de cartera que cobra costes de operaciones no
ejecutadas: todos terminaron su ejecución con éxito aparente y escribieron sus artefactos.

La cuarta es sobre estimadores. Un diagnóstico por deciles basado en la media sugirió durante un
tiempo que la ordenación estaba invertida en todo el histórico; con la mediana, la conclusión se
invertía. Antes de rediseñar un sistema conviene comprobar que el estadístico empleado para
diagnosticarlo sea robusto a la distribución que realmente tienen los datos.

La quinta, y la más incómoda, es que conviene registrar las decisiones tomadas en contra de la
recomendación técnica, con ambos lados del argumento. No mejoran el resultado, pero son
precisamente las que un lector externo no podría reconstruir desde el código.
